\chapter{Downtime}
\label{ch:downtime}

\begin{tcolorbox}[colback=blue!5!white,colframe=blue!75!black,title=When Does Downtime Happen?]
Downtime occurs between adventures, when the party returns to a safe haven (a village, a guild hall, a camp that won't be attacked). Usually you take several days or weeks ``off‑screen'' while the world advances. Your GM will tell you when downtime begins.
\end{tcolorbox}

\section{What You Can Do During Downtime}
Choose any number of the following activities. Each one takes a specific amount of time (usually a day or a week). You may split your downtime between multiple activities, but you cannot be in two places at once.

\subsection{Recovery}
\begin{itemize}
  \item \textbf{Clear Fatigue:} A short rest (1 hour) removes 2 Fatigue. A full night's rest (8 hours) removes all Fatigue.
  \item \textbf{Heal Harm:} Without magic, you restore 1 level of Harm per week of complete rest. A successful \textbf{Wits + Medicine} roll (DV 3) at the start of the week reduces the time to 3 days per Harm level.
  \item \textbf{Remove Conditions:} Use medicine, magic, or ritual to purge a Condition (e.g., Poisoned, Cursed). This usually requires a \textbf{Spirit + Lore} roll (DV 3--5) and a full day.
\end{itemize}

\subsection{Training}
You may spend XP to improve Attributes, Skills, or Talents. Each improvement requires training time (see table). You can train multiple things in parallel if you have enough downtime days.

\begin{center}
\begin{tabular}{lcc}
\toprule
\textbf{Improvement} & \textbf{XP Cost} & \textbf{Training Time} \\
\midrule
Attribute increase & new rating \(\times\)3 XP & new rating in days \\
Skill increase & new level \(\times\)2 XP & new level in days \\
New Talent & listed XP & 3--10 days (GM discretion) \\
\bottomrule
\end{tabular}
\end{center}

\paragraph{Rush Rule:} If you want to complete training faster, you may attempt a \textbf{Rush}. The GM starts a \textbf{Risk Timer [4]}. Each downtime segment you roll \textbf{Spirit + Resolve} (DV 3) to push yourself. For each failure (or 1 on the die), the timer ticks. If the timer fills before training is complete, the new ability has a flaw (temporary –1 die, or a narrative complication). The GM decides the flaw.

\subsection{Asset \& Follower Upkeep}
Assets (safehouses, workshops) and Followers (allies, pets) must be maintained each downtime. Choose one method per asset or follower:

\begin{itemize}
  \item \textbf{Efficient (higher XP, less time):} Pay XP equal to the acquisition cost divided by 3 (rounded up, minimum 1). No scene required.
  \item \textbf{Intensive (lower XP, more time):} Pay 1 XP and spend a scene describing how you maintain the asset or relationship (repairing the roof, training with your follower). This counts as a downtime activity.
\end{itemize}
If you pay neither, the asset/follower becomes \textbf{Neglected} (no benefit until you pay retroactive upkeep). Two consecutive neglected periods make it \textbf{Compromised} (must be repaired/recovered through a quest).

\subsection{Clearing Obligation (to a Patron or Terrestrial Patron)}
You may clear \textbf{Obligation} segments (see \ref{ch:world-powers-obligation}) by performing an \textbf{Act of Service}. Describe the act in detail (intricate description) – it must align with the Patron's themes. For example:
\begin{itemize}
  \item \textbf{Ikasha:} Burn an old cloak, symbolising leaving behind a past self.
  \item \textbf{The Traveler:} Guide a lost traveler to safety.
  \item \textbf{Oath of Flame \& Light:} Swear a public vow and keep it for a day.
\end{itemize}
Each such act clears \textbf{1 segment} of Obligation (off‑screen). You may also use the optional \textbf{Focused Devotion} talent (3--5 XP) to clear 1 segment on‑screen by spending 1 Boon (once per scene).

\subsection{Crafting}
\label{subsec:downtime-crafting}
See \ref{ch:gear} for full rules (the section on magic items and crafting). In brief:
\begin{itemize}
  \item \textbf{Mundane items:} No roll. You can make tools, clothing, etc. if you have Craft skill.
  \item \textbf{Consumable magic items (potions, scrolls):} Takes one downtime action, costs 1 XP (or 2 XP for powerful ones). Roll \textbf{Wits + Arcana} or \textbf{Wits + Craft} (DV 3). On failure, components are wasted.
  \item \textbf{Permanent magic items:} Requires a Downtime Project Timer (6--12 segments), rare components (obtained through quests), and the XP cost of the item (see \ref{ch:gear}). You roll each downtime to advance the timer.
\end{itemize}

\subsection{Research}
You may spend downtime investigating a mystery, learning lore, or preparing for an upcoming challenge.

\begin{itemize}
  \item \textbf{Basic research:} Roll \textbf{Wits + Lore} (DV 3). On success, you learn one useful fact (GM discretion). On partial, you learn something, but it's incomplete or costs extra time.
  \item \textbf{Extended research:} Use a Progress Timer [4--8] for complex topics (e.g., uncovering a noble's secret). Each downtime you can roll to advance the timer.
\end{itemize}

\subsection{Carousing \& Networking}
Spend time socialising to gain rumors, contacts, or temporary favors.

\begin{itemize}
  \item \textbf{Rumors:} Roll \textbf{Presence + Sway} (DV 3). On success, learn one rumor (true or false – the GM decides). On partial, the rumor is vague or costs you a drink (1 Fatigue).
  \item \textbf{Gain a temporary contact:} Spend 1 XP to establish a minor contact who can grant +1 die on one relevant roll in the next session. The contact fades after use.
\end{itemize}

\subsection{Work \& Income}
In a currency‑free game, ``income'' is abstract. You may:
\begin{itemize}
  \item \textbf{Perform a job:} Spend a downtime week working. You earn enough to cover your expenses and have a small amount of barter goods (narrative permission for minor bribes or gear). No roll needed unless you want extra profit – then roll \textbf{Presence + Sway} or \textbf{Wits + Craft} (DV 4). On success, you gain 1 XP or a minor consumable item.
\end{itemize}

\subsection{Special Downtime Projects}
Your character may have a unique long‑term goal (e.g., building a guild, writing a grimoire, taming a monster). Work with your GM to create a \textbf{Downtime Project Timer } (6--12 segments). Each downtime you can describe your progress and roll an appropriate skill to advance the timer. When it fills, the project is complete.

\section{Off‑Screen Character Switching (Echo Characters)}
If your table uses the optional multi‑character rule (see \ref{ch:enhanced-play}), you may send one character off‑screen while you play another. The off‑screen character maintains a \textbf{Downtime Log}:
\begin{itemize}
  \item \textbf{Resource Action:} How they maintain assets or relationships.
  \item \textbf{Project Action:} One goal or project advanced.
  \item \textbf{Event Response:} How they react to unseen events.
\end{itemize}
The GM records consequences. When you return to that character, you gain 2 Boons and roll \textbf{Wits + Spirit} (DV 3) to re‑establish presence.

\section{GM's Role (Reminder)}
Your GM determines how much downtime you have between adventures (usually a few days to several weeks). They also set DVs for downtime rolls and adjudicate any complications. Downtime is not a ``free'' rest – the world may advance timers (faction agendas, rival plans) while you rest. Ask your GM what has changed when you return.

\begin{tcolorbox}[colback=green!5!white,colframe=green!70!black,title=The Gray Wanderer's Advice]
``Downtime is not a pause. It is a negotiation. You mend your boots, but the road mends itself. When you step back onto it, it will be different – and so will you.''
\end{tcolorbox}